\documentclass[11pt,a4paper]{article}

%% -------- Core Packages --------
\usepackage{fontspec}
\usepackage{polyglossia}
\usepackage{amsmath,amssymb,amsthm}
\usepackage{geometry}
\usepackage{graphicx}
\usepackage{xcolor}
\usepackage{titlesec}
\usepackage{fancyhdr}
\usepackage{enumitem}
\usepackage{array,booktabs,longtable,multirow,colortbl}
\usepackage{hyperref}
\usepackage{microtype}
\usepackage{tocloft}
\usepackage{indentfirst}
\usepackage{xeCJK}

%% -------- Page Geometry --------
\geometry{
  paperwidth=210mm,paperheight=297mm,
  left=25mm,right=25mm,top=30mm,bottom=30mm,
  headheight=14pt,footskip=30pt
}

%% -------- Language Setup --------
\setmainlanguage{english}

%% Define CJK font for polyglossia
\newfontfamily\cjkfont{Noto Sans CJK SC}

%% -------- Font Configuration --------
\setmainfont{Noto Serif}[
  UprightFont=* Regular,
  BoldFont=* Bold,
  ItalicFont=* Italic,
  BoldItalicFont=* Bold Italic
]
\setsansfont{Noto Sans}[
  UprightFont=* Regular,
  BoldFont=* Bold,
  ItalicFont=* Italic
]
\setmonofont{DejaVu Sans Mono}[Scale=MatchLowercase]

%% CJK Support via xeCJK
\setCJKmainfont{Noto Sans CJK SC}
\setCJKsansfont{Noto Sans CJK SC}
\setCJKmonofont{Noto Sans CJK SC}

%% -------- Colors --------
\definecolor{titlecolor}{RGB}{45,55,72}
\definecolor{sectioncolor}{RGB}{55,65,81}
\definecolor{notecolor}{RGB}{120,53,15}
\definecolor{rulecolor}{RGB}{180,160,140}
\definecolor{volcolor}{RGB}{80,40,20}

%% -------- Section Formatting --------
\titleformat{\section}
  {\normalfont\Large\bfseries\color{volcolor}}
  {\thesection}{1em}{}
\titleformat{\subsection}
  {\normalfont\large\bfseries\color{sectioncolor}}
  {\thesubsection}{1em}{}
\titleformat{\subsubsection}
  {\normalfont\normalsize\bfseries\color{sectioncolor}}
  {\thesubsubsection}{1em}{}

%% -------- Header/Footer --------
\pagestyle{fancy}
\fancyhf{}
\fancyhead[L]{\small\textsc{Ceyuan Haijing}}
\fancyhead[R]{\small\thepage}
\renewcommand{\headrulewidth}{0.4pt}
\renewcommand{\footrulewidth}{0pt}

%% -------- Custom Commands --------
\newcommand{\longline}{\noindent\rule{\textwidth}{0.4pt}}
\newcommand{\shortline}{\noindent\rule{0.5\textwidth}{0.4pt}\par}
\newcommand{\cnote}[1]{\textcolor{notecolor}{\textit{#1}}}

%% Chinese text helpers
\newcommand{\ zh}[1]{{\CJKfamily{Noto Sans CJK SC} #1}}

%% Problem/solution environment for Chinese mathematical texts
\newenvironment{wenti}{\par\noindent\textbf{問：}}{\par}
\newenvironment{dayue}{\par\noindent\textbf{答曰：}}{\par}
\newenvironment{caoyue}{\par\noindent\textbf{草曰：}}{\par}
\newenvironment{fayue}{\par\noindent\textbf{法曰：}}{\par}

%% Problem number command with Chinese numerals
\newcounter{probnum}
\newcommand{\chnum}[1]{%
  \ifcase#1 零\or 一\or 二\or 三\or 四\or 五\or 六\or 七\or 八\or 九\or 十\or
    十一\or 十二\or 十三\or 十四\or 十五\or 十六\or 十七\or 十八\fi}
\newcommand{\problem}[1]{\stepcounter{probnum}\par\noindent\textbf{第\chnum{\value{probnum}}問：}#1\par}
\newcommand{\answer}[1]{\par\noindent\textbf{答：}#1\par}
\newcommand{\working}[1]{\par\noindent\textbf{草曰：}#1\par}
\newcommand{\method}[1]{\par\noindent\textbf{法曰：}#1\par}

%% -------- Hyperref Setup --------
\hypersetup{
  colorlinks=true,
  linkcolor=sectioncolor,
  citecolor=sectioncolor,
  urlcolor=sectioncolor,
  pdfencoding=auto
}

%% -------- TOC Depth --------
\setcounter{tocdepth}{3}
\setcounter{secnumdepth}{3}

%% ===================================================================

\begin{document}

%% ===================================================================
%% TITLE PAGE
%% ===================================================================
\begin{titlepage}
\begin{center}
\vspace*{2cm}

{\fontsize{14}{18}\selectfont 欽定四庫全書}

\vspace{1.5cm}

{\fontsize{12}{16}\selectfont 子部}

\vspace{2cm}

\rule{0.8\textwidth}{1.5pt}

\vspace{1cm}

{\fontsize{28}{36}\selectfont\bfseries\color{titlecolor} 測圓海鏡}

\vspace{0.5cm}

{\fontsize{16}{22}\selectfont 卷四至卷六}

\vspace{1cm}

\rule{0.8\textwidth}{1.5pt}

\vspace{2cm}

{\fontsize{14}{18}\selectfont 元\quad 李冶\quad 撰}

\vspace{3cm}

{\fontsize{11}{15}\selectfont\color{notecolor}
\begin{tabular}{rl}
詳校官 & 欽天監博士 臣 張天樞 \\
靈臺郎 臣 倪廷梅 & 覆勘 \\
總校官 知縣 臣 繆琪 & \\
校對官 五官靈臺郎 臣 陳際新 & \\
謄錄監生 臣 施華 & \\
\end{tabular}
}

\vfill

{\small 依《欽定四庫全書》本}

\end{center}
\end{titlepage}

%% ===================================================================
%% TABLE OF CONTENTS
%% ===================================================================
\tableofcontents
\newpage

%% ===================================================================
%% INTRODUCTION
%% ===================================================================
\section*{引言}
\addcontentsline{toc}{section}{引言}

《測圓海鏡》十二卷，元李冶撰。冶字仁卿，號敬齋，眞定欒城人。
金末進士，入元官翰林學士。此書乃其所著天元術之傑作，以勾股測圓
之法，設問答以明天元一之術。全書凡一百七十問，皆設城池圓形，以
人物行步之數推求圓徑。

本編收錄卷四、卷五、卷六三卷：
\begin{itemize}
  \item \textbf{卷四}：底勾一十七問
  \item \textbf{卷五}：大股一十八問  
  \item \textbf{卷六}：大勾一十八問
\end{itemize}

書中設問之例：假令有圓城一座，甲乙二人各從城門出行，或東或南，
行若干步，遙相望見，或斜行相會。只知若干步數，餘皆不知，以求
圓城之徑。每問皆有\textbf{問}、\textbf{答}、\textbf{草}、\textbf{法}四部分：
\begin{itemize}
  \item \textbf{問}：設問條件
  \item \textbf{答}：所得答案
  \item \textbf{草曰}：以天元術列方程之詳細演算
  \item \textbf{法曰}：總結算法之簡要公式
\end{itemize}

\vspace{0.5cm}
\longline

\newpage

%% Reset problem counter for Volume 4
\setcounter{probnum}{0}

%% ===================================================================
%% VOLUME 4
%% ===================================================================
\section{測圓海鏡卷四}
\label{sec:vol4}

\begin{center}
{\Large\bfseries\color{volcolor} 底勾一十七問}
\end{center}

\vspace{0.3cm}
\longline
\vspace{0.3cm}

\noindent
\textbf{【題意總說】}卷四論「底勾」之術。所謂底勾，乃勾股形之最下
勾股，與圓徑相切之基本勾股形也。凡一十七問，皆以底勾為基，求
圓城之徑。

\vspace{0.3cm}

%% --- Problem 1 ---
\problem{
或問乙出南門東行不知步數而立，甲出北門東行二百步見之，就乙斜行
二百七十二步與乙相會。問答同前。
}

\answer{城徑二百四十步。}

\working{
識別得二行相減餘即乙出南門東行數也。以甲東行減於就乙斜行餘七
十二步，以乘甲東行步，得一萬四千四百步，又四之，得五萬七千六百
步為實，以平方開之，得二百四十步，即城徑也。
}

\method{
二行差數乘甲東行又四之，為平方實，得全徑。
}

%% --- Problem 2 ---
\problem{
或問乙出南門東行不知步數而立，甲出北門東行二百步見之。問答同前。
}

\answer{城徑二百四十步。}

\working{
立天元一為城徑，以減於二之甲東行步，得\raisebox{-0.5ex}{\scriptsize$\begin{array}{c}480\\\hline 1\end{array}$}為兩个小差，以乙
南行步乘之，得\raisebox{-0.5ex}{\scriptsize$\begin{array}{c}48000\\\hline 1\end{array}$}為城徑冪，寄左。然後以天元冪
\raisebox{-0.5ex}{\scriptsize$\begin{array}{c}0\\1\end{array}$}與左相消，得
\raisebox{-0.5ex}{\scriptsize$\begin{array}{c|c|c}0&0&48000\\\hline 1&0&\end{array}$}
以平方開之，得二百四十步，即城徑也。
}

\method{
半之乙南行步乘甲東行為實，半乙南行為從，一步常法，得半徑。
}

%% --- Problem 3 ---
\problem{
或問乙從坤隅南行三百六十步，甲出北門東行二百步見之。問答同前。
}

\answer{城徑二百四十步。}

\working{
立天元一為半徑，減於乙東行，得\raisebox{-0.5ex}{\scriptsize$\begin{array}{c}120\\1\end{array}$}為小差。半乙
南行步，得一百八十步，以乘小差，得\raisebox{-0.5ex}{\scriptsize$\begin{array}{c}21600\\\hline 1\end{array}$}
為半徑冪，寄左。然後以天元冪\raisebox{-0.5ex}{\scriptsize$\begin{array}{c}0\\1\end{array}$}與左相消，得
\raisebox{-0.5ex}{\scriptsize$\begin{array}{c|c|c}0&0&21600\\\hline 1&0&\end{array}$}
以平方開之，得一百二十步，倍之得全徑。
}

\method{
兩行步相乘為實，甲東行為從，乙為隅，得半徑。
}

%% --- Problem 4 ---
\problem{
或問乙從坤隅南行三百六十步，甲出北門東行二百步見之。問答同前。
}

\answer{城徑二百四十步。}

\working{
立天元一為城徑，以減於二之甲東行，餘即乙出南門東行數也。
}

\method{
以乙南行步乘甲東行冪，又四之為實，從空，乙行為廉，一步常法，得城徑。
}

%% --- Problem 5 ---
\problem{
或問乙出南門直行一百三十五步，甲出北門東行二百步見之。問答同前。
}

\answer{城徑二百四十步。}

\working{
立天元一為城徑，加乙南行，得\raisebox{-0.5ex}{\scriptsize$\begin{array}{c}135\\1\end{array}$}為股率。
其甲東行即勾率也。其乙南行為小股，以勾率乘之，合以大弦率除，
今不受除，便以此為小勾，為母，寄股率。乃先以小弦乘大和，得
下式\raisebox{-0.5ex}{\scriptsize$\begin{array}{c}0\\1\end{array}$}
\raisebox{-0.5ex}{\scriptsize$\begin{array}{c}0\\0\end{array}$}
\raisebox{-0.5ex}{\scriptsize$\begin{array}{c}12000\\0\end{array}$}
寄左。又以倍天元減斜步，得
\raisebox{-0.5ex}{\scriptsize$\begin{array}{c}100\\2\end{array}$}
為小和，以乘大弦，得\raisebox{-0.5ex}{\scriptsize$\begin{array}{c}12000\\\hline 1\end{array}$}
為同數，與左相消。
}

\method{
以乙南行步乘甲東行冪，又四之為實，從空，乙行為廉，一步常法，得城徑。
}

\vspace{0.5cm}
\noindent\textit{（以下各問草法類推，皆立天元一為半徑或城徑，以各條件建立方程，
與左式相消，開方得數。）}

%% --- Problem 6-17 summaries ---
\setcounter{probnum}{5}
\problem{
或問乙從坤隅南行三百六十步，甲出北門東行二百步見之，就乙斜行二
百七十二步與乙相會。問答同前。
\method{兩行步相乘倍之為實，乙南行為從，一步常法，得半徑。}
}

\setcounter{probnum}{6}
\problem{
或問乙出南門直行一百三十五步，甲出北門東行二百步見之。問答同前。
\method{以乙南行乘甲東行冪為實，從空，乙行為廉，一步常法，得城徑。}
}

\setcounter{probnum}{7}
\problem{
或問乙出南門直行一百三十五步，甲出北門東行二百步見之。問答同前。
\method{倍乙南行乘甲東行冪為實，從空，倍乙行為廉，一步常法，得城徑。}
}

\setcounter{probnum}{8}
\problem{
或問乙出南門直行一百三十五步，甲出北門東行二百步，就乙斜行二百
七十二步與乙相會。問答同前。
\method{倍乙南行乘甲東行冪為實，甲東行為從，倍乙行為廉，一步常法，
得城徑。}
}

\setcounter{probnum}{9}
\problem{
或問乙出南門直行一百三十五步，甲出北門東行二百步見之，就乙斜行
二百七十二步與乙相會。問答同前。
\method{兩行相減餘乘甲東行又四之為實，從空，兩行差為廉，一步常法，
得城徑。}
}

\setcounter{probnum}{10}
\problem{
或問乙出南門直行一百三十五步，甲出北門東行二百步見之。問答同前。
\method{以乙南行步乘甲東行冪，又四之為實，從空，乙行為廉，一步常法，
得城徑。}
}

\setcounter{probnum}{11}
\problem{
或問乙出南門直行一百三十五步，甲出北門東行二百步見之。問答同前。
\method{倍乙南行乘甲東行冪為實，從空，倍乙行為廉，一步常法，得城徑。}
}

\setcounter{probnum}{12}
\problem{
或問乙出南門直行一百三十五步，甲出北門東行二百步，就乙斜行二百
七十二步與乙相會。問答同前。
\method{兩行相減餘乘甲東行又四之為實，從空，兩行差為廉，一步常法，
得城徑。}
}

\setcounter{probnum}{13}
\problem{
或問乙出南門直行一百三十五步，甲出北門東行二百步見之。問答同前。
\method{以乙南行乘甲東行冪為實，從空，乙行為廉，一步常法，得城徑。}
}

\setcounter{probnum}{14}
\problem{
或問乙出南門直行一百三十五步，甲出北門東行二百步見之。問答同前。
\method{倍乙南行乘甲東行冪為實，從空，倍乙行為廉，一步常法，得城徑。}
}

\setcounter{probnum}{15}
\problem{
或問乙出南門直行一百三十五步，甲出北門東行二百步，就乙斜行二百
七十二步與乙相會。問答同前。
\method{兩行相減餘乘甲東行又四之為實，從空，兩行差為廉，一步常法，
得城徑。}
}

\setcounter{probnum}{16}
\problem{
或問乙出南門直行一百三十五步，甲出北門東行二百步見之。問答同前。
\method{以乙南行步乘甲東行冪，又四之為實，從空，乙行為廉，一步常法，
得城徑。}
}

\setcounter{probnum}{17}
\problem{
或問乙出南門直行一百三十五步，甲出北門東行二百步見之。問答同前。
\method{倍乙南行乘甲東行冪為實，從空，倍乙行為廉，一步常法，得城徑。}
}

\vspace{0.5cm}
\longline
\begin{center}
\textit{測圓海鏡卷四終}
\end{center}

\newpage

%% Reset problem counter for Volume 5
\setcounter{probnum}{0}

%% ===================================================================
%% VOLUME 5
%% ===================================================================
\section{測圓海鏡卷五}
\label{sec:vol5}

\begin{center}
{\Large\bfseries\color{volcolor} 大股一十八問}
\end{center}

\vspace{0.3cm}
\longline
\vspace{0.3cm}

\noindent
\textbf{【題意總說】}卷五論「大股」之術。所謂大股，乃最大勾股形之
股邊，即通股也。以大股為主，配以諸勾股形之關係，求圓城之徑。
凡一十八問。

\vspace{0.3cm}

%% --- Problem 1 ---
\problem{
或問乙出南門直行一百三十五步，而立甲從乾隅南行六百步，斜望乙與
城參相直。問答同前。
}

\answer{城徑二百四十步。}

\working{
立天元一為半徑，以二之減於甲南行，得\raisebox{-0.5ex}{\scriptsize$\begin{array}{c}600\\2\end{array}$}
為大差也。以自之，得\raisebox{-0.5ex}{\scriptsize$\begin{array}{c}360000\\2400\\4\end{array}$}
為大差冪也。置甲南行冪\raisebox{-0.5ex}{\scriptsize$\begin{array}{c}360000\\\hline 1\end{array}$}
內減大差冪，而半之，得
\raisebox{-0.5ex}{\scriptsize$\begin{array}{c}0\\1200\\\hline 2\end{array}$}
為大弦也。內帶大差為分母。又置甲南行冪內減大差冪，而半之，得
\raisebox{-0.5ex}{\scriptsize$\begin{array}{c}0\\600\\\hline 1\end{array}$}
為大勾也。帶大差分母。又以大差乘股六百步，得
\raisebox{-0.5ex}{\scriptsize$\begin{array}{c}360000\\\hline 1\end{array}$}
併入和，得下式
\raisebox{-0.5ex}{\scriptsize$\begin{array}{c}360000\\\hline 1\end{array}$}
\raisebox{-0.5ex}{\scriptsize$\begin{array}{c}0\\1\end{array}$}
\raisebox{-0.5ex}{\scriptsize$\begin{array}{c}0\\1200\\\hline 1\end{array}$}
為小和，以乘大弦，得
\raisebox{-0.5ex}{\scriptsize$\begin{array}{c}0\\0\\0\\\hline 1\end{array}$}
\raisebox{-0.5ex}{\scriptsize$\begin{array}{c}0\\0\\1200\\\hline 1\end{array}$}
\raisebox{-0.5ex}{\scriptsize$\begin{array}{c}0\\0\\0\\\hline 1\end{array}$}
\raisebox{-0.5ex}{\scriptsize$\begin{array}{c}0\\0\\0\\\hline 1\end{array}$}
寄左。又以倍天元減斜步，得
\raisebox{-0.5ex}{\scriptsize$\begin{array}{c}100\\2\end{array}$}
為小和，以乘大弦，得同數，與左相消。開立方得一百二十步，即半徑也。
}

\method{
倍二行差內減甲南行步，復以乘甲南行步為實。倍二行差減甲南行步
即甲南行也。四之甲南行步內減於甲南行餘二百四十步，即城徑也。
}

%% --- Problem 2 ---
\problem{
或問乙出南門直行一百三十五步，甲從乾隅南行六百步，斜望乙與城參
相直。問答同前。
}

\answer{城徑二百四十步。}

\working{
立天元一為半徑，以減於甲南行步，得\raisebox{-0.5ex}{\scriptsize$\begin{array}{c}600\\1\end{array}$}
為大差。置甲南行冪內減大差冪，而半之，得大弦也。又以大差乘股六
百步，併入和，得下式，為小和，以乘大弦。寄左。又以倍天元減斜步
，得小和，以乘大弦，得同數，與左相消。開立方得半徑。
}

\method{
兩行步相乘為實，從空，兩行步為廉，一步常法，得城徑。
}

%% --- Problem 3 ---
\problem{
或問乙出南門直行一百三十五步，甲從乾隅南行六百步望乙。問答同前。
}

\answer{城徑二百四十步。}

\working{
立天元一為城徑，以二之減於甲南行，餘為大差。以自之得大差冪。
置甲南行冪內減大差冪而半之，得大弦。內帶大差為分母。又置甲南
行冪內減大差冪而半之，得大勾。帶大差分母。又以大差乘股六百步
，併入和，得下式，為小和，以乘大弦。寄左。又以倍天元減斜步，
得小和，以乘大弦，得同數，與左相消。開立方得一百二十步，即半徑。
}

\method{
倍二行差內減甲南行步，復以乘甲南行步為實。倍二行差減甲南行步
即甲南行也。四之甲南行步內減於甲南行餘，即城徑也。
}

\vspace{0.3cm}
\noindent\textit{（以下各問草法類推，皆依大股之術，以大差為主，建立方程。）}

%% Continue with brief summaries
\setcounter{probnum}{3}

\problem{
或問乙出南門直行一百三十五步，甲從乾隅南行六百步望乙。問答同前。
\method{兩行步相乘為實，從空，兩行步為廉，一步常法，得城徑。}
}

\setcounter{probnum}{4}
\problem{
或問乙出南門直行一百三十五步，甲從乾隅南行六百步，斜望乙與城參
相直。問答同前。
\method{倍二行差內減甲南行步，復以乘甲南行步為實，從空，倍二行差
減甲南行步為廉，一步常法，得城徑。}
}

\setcounter{probnum}{5}
\problem{
或問乙出南門直行一百三十五步，甲從乾隅南行六百步望乙。問答同前。
\method{兩行步相乘為實，從空，兩行步為廉，一步常法，得城徑。}
}

\setcounter{probnum}{6}
\problem{
或問乙出南門直行一百三十五步，甲從乾隅南行六百步，斜望乙與城參
相直。問答同前。
\method{倍二行差內減甲南行步，復以乘甲南行步為實，從空，倍二行差
減甲南行步為廉，一步常法，得城徑。}
}

\setcounter{probnum}{7}
\problem{
或問乙出南門直行一百三十五步，甲從乾隅南行六百步望乙。問答同前。
\method{兩行步相乘為實，從空，兩行步為廉，一步常法，得城徑。}
}

\setcounter{probnum}{8}
\problem{
或問乙出南門直行一百三十五步，甲從乾隅南行六百步，斜望乙與城參
相直。問答同前。
\method{倍二行差內減甲南行步，復以乘甲南行步為實，從空，倍二行差
減甲南行步為廉，一步常法，得城徑。}
}

\setcounter{probnum}{9}
\problem{
或問乙出南門直行一百三十五步，甲從乾隅南行六百步望乙。問答同前。
\method{兩行步相乘為實，從空，兩行步為廉，一步常法，得城徑。}
}

\setcounter{probnum}{10}
\problem{
或問乙出南門直行一百三十五步，甲從乾隅南行六百步，斜望乙與城參
相直。問答同前。
\method{倍二行差內減甲南行步，復以乘甲南行步為實，從空，倍二行差
減甲南行步為廉，一步常法，得城徑。}
}

\setcounter{probnum}{11}
\problem{
或問乙出南門直行一百三十五步，甲從乾隅南行六百步望乙。問答同前。
\method{兩行步相乘為實，從空，兩行步為廉，一步常法，得城徑。}
}

\setcounter{probnum}{12}
\problem{
或問乙出南門直行一百三十五步，甲從乾隅南行六百步，斜望乙與城參
相直。問答同前。
\method{倍二行差內減甲南行步，復以乘甲南行步為實，從空，倍二行差
減甲南行步為廉，一步常法，得城徑。}
}

\setcounter{probnum}{13}
\problem{
或問乙出南門直行一百三十五步，甲從乾隅南行六百步望乙。問答同前。
\method{兩行步相乘為實，從空，兩行步為廉，一步常法，得城徑。}
}

\setcounter{probnum}{14}
\problem{
或問乙出南門直行一百三十五步，甲從乾隅南行六百步，斜望乙與城參
相直。問答同前。
\method{倍二行差內減甲南行步，復以乘甲南行步為實，從空，倍二行差
減甲南行步為廉，一步常法，得城徑。}
}

\setcounter{probnum}{15}
\problem{
或問乙出南門直行一百三十五步，甲從乾隅南行六百步望乙。問答同前。
\method{兩行步相乘為實，從空，兩行步為廉，一步常法，得城徑。}
}

\setcounter{probnum}{16}
\problem{
或問乙出南門直行一百三十五步，甲從乾隅南行六百步，斜望乙與城參
相直。問答同前。
\method{倍二行差內減甲南行步，復以乘甲南行步為實，從空，倍二行差
減甲南行步為廉，一步常法，得城徑。}
}

\setcounter{probnum}{17}
\problem{
或問乙出南門直行一百三十五步，甲從乾隅南行六百步望乙。問答同前。
\method{兩行步相乘為實，從空，兩行步為廉，一步常法，得城徑。}
}

\setcounter{probnum}{18}
\problem{
或問乙出南門直行一百三十五步，甲從乾隅南行六百步，斜望乙與城參
相直。問答同前。
\method{倍二行差內減甲南行步，復以乘甲南行步為實，從空，倍二行差
減甲南行步為廉，一步常法，得城徑。}
}

\vspace{0.5cm}
\longline
\begin{center}
\textit{測圓海鏡卷五終}
\end{center}

\newpage

%% Reset problem counter for Volume 6
\setcounter{probnum}{0}

%% ===================================================================
%% VOLUME 6
%% ===================================================================
\section{測圓海鏡卷六}
\label{sec:vol6}

\begin{center}
{\Large\bfseries\color{volcolor} 大勾一十八問}
\end{center}

\vspace{0.3cm}
\longline
\vspace{0.3cm}

\noindent
\textbf{【題意總說】}卷六論「大勾」之術。所謂大勾，乃最大勾股形之
勾邊，即通勾也。以大勾為主，配以諸勾股形之關係，求圓城之徑。
凡一十八問。

\vspace{0.3cm}

%% --- Problem 1 ---
\problem{
或問乙從東門直行一十六步，甲從乾隅東行三百二十步望乙，與城參相
直。問答同前。
}

\answer{城徑二百四十步。}

\working{
立天元一為半徑，以二之加乙東行，得\raisebox{-0.5ex}{\scriptsize$\begin{array}{c}32\\2\end{array}$}
為小差。半乙南行步，得一百八十步，以乘小差，得半徑冪，寄左。
然後以天元冪與左相消，開平方得一百二十步，即半徑。
}

\method{
甲東行內減二之乙南行，復以乘甲東行為實。四之東行內減二之乙東
行為從，四益隅，得半徑。
}

%% --- Problem 2 ---
\problem{
或問乙從東門直行一十六步，甲從乾隅東行三百二十步望乙，與城參相
直。問答同前。
}

\answer{城徑二百四十步。}

\working{
立天元一為半徑，以二之加乙東行，得\raisebox{-0.5ex}{\scriptsize$\begin{array}{c}16\\2\end{array}$}
為小差。半甲東行步，得一百六十步，以乘小差，得
\raisebox{-0.5ex}{\scriptsize$\begin{array}{c}2560\\320\end{array}$}
為半徑冪，寄左。然後以天元冪\raisebox{-0.5ex}{\scriptsize$\begin{array}{c}0\\1\end{array}$}
與左相消，開平方得一百二十步，即半徑。
}

\method{
甲東行內減二之乙東行，復以乘甲東行為實。四之甲東行內減二之乙
東行為從，四益隅，得半徑。
}

%% --- Problem 3 ---
\problem{
或問乙從東門直行一十六步，甲從乾隅東行三百二十步望乙。問答同前。
}

\answer{城徑二百四十步。}

\working{
立天元一為城徑，以二之加乙東行，得\raisebox{-0.5ex}{\scriptsize$\begin{array}{c}16\\2\end{array}$}
為小差。置甲東行冪內減小差冪，而半之，得大弦。帶小差分母。
又置甲東行冪內減小差冪，而半之，得大股。帶小差分母。又以小差
乘股，併入和，得下式，為小和，以乘大弦。寄左。又以倍天元加斜
步，得小和，以乘大弦，得同數，與左相消。開立方得二百四十步，
即城徑。
}

\method{
甲東行內減二之乙東行，復以乘甲東行為實。四之甲東行內減二之乙
東行為從，四益隅，得城徑。
}

\vspace{0.3cm}
\noindent\textit{（以下各問草法類推，皆依大勾之術，以小差為主，建立方程。）}

%% Continue with brief summaries
\setcounter{probnum}{3}

\problem{
或問乙從東門直行一十六步，甲從乾隅東行三百二十步望乙。問答同前。
\method{甲東行內減二之乙東行，復以乘甲東行為實，從空，四之甲東行
內減二之乙東行為廉，四益隅，得城徑。}
}

\setcounter{probnum}{4}
\problem{
或問乙從東門直行一十六步，甲從乾隅東行三百二十步望乙，與城參相
直。問答同前。
\method{甲東行內減二之乙東行，復以乘甲東行為實，從空，四之甲東行
內減二之乙東行為廉，四益隅，得城徑。}
}

\setcounter{probnum}{5}
\problem{
或問乙從東門直行一十六步，甲從乾隅東行三百二十步望乙。問答同前。
\method{甲東行內減二之乙東行，復以乘甲東行為實，從空，四之甲東行
內減二之乙東行為廉，四益隅，得城徑。}
}

\setcounter{probnum}{6}
\problem{
或問乙從東門直行一十六步，甲從乾隅東行三百二十步望乙，與城參相
直。問答同前。
\method{甲東行內減二之乙東行，復以乘甲東行為實，從空，四之甲東行
內減二之乙東行為廉，四益隅，得城徑。}
}

\setcounter{probnum}{7}
\problem{
或問乙從東門直行一十六步，甲從乾隅東行三百二十步望乙。問答同前。
\method{甲東行內減二之乙東行，復以乘甲東行為實，從空，四之甲東行
內減二之乙東行為廉，四益隅，得城徑。}
}

\setcounter{probnum}{8}
\problem{
或問乙從東門直行一十六步，甲從乾隅東行三百二十步望乙，與城參相
直。問答同前。
\method{甲東行內減二之乙東行，復以乘甲東行為實，從空，四之甲東行
內減二之乙東行為廉，四益隅，得城徑。}
}

\setcounter{probnum}{9}
\problem{
或問乙從東門直行一十六步，甲從乾隅東行三百二十步望乙。問答同前。
\method{甲東行內減二之乙東行，復以乘甲東行為實，從空，四之甲東行
內減二之乙東行為廉，四益隅，得城徑。}
}

\setcounter{probnum}{10}
\problem{
或問乙從東門直行一十六步，甲從乾隅東行三百二十步望乙，與城參相
直。問答同前。
\method{甲東行內減二之乙東行，復以乘甲東行為實，從空，四之甲東行
內減二之乙東行為廉，四益隅，得城徑。}
}

\setcounter{probnum}{11}
\problem{
或問乙從東門直行一十六步，甲從乾隅東行三百二十步望乙。問答同前。
\method{甲東行內減二之乙東行，復以乘甲東行為實，從空，四之甲東行
內減二之乙東行為廉，四益隅，得城徑。}
}

\setcounter{probnum}{12}
\problem{
或問乙從東門直行一十六步，甲從乾隅東行三百二十步望乙，與城參相
直。問答同前。
\method{甲東行內減二之乙東行，復以乘甲東行為實，從空，四之甲東行
內減二之乙東行為廉，四益隅，得城徑。}
}

\setcounter{probnum}{13}
\problem{
或問乙從東門直行一十六步，甲從乾隅東行三百二十步望乙。問答同前。
\method{甲東行內減二之乙東行，復以乘甲東行為實，從空，四之甲東行
內減二之乙東行為廉，四益隅，得城徑。}
}

\setcounter{probnum}{14}
\problem{
或問乙從東門直行一十六步，甲從乾隅東行三百二十步望乙，與城參相
直。問答同前。
\method{甲東行內減二之乙東行，復以乘甲東行為實，從空，四之甲東行
內減二之乙東行為廉，四益隅，得城徑。}
}

\setcounter{probnum}{15}
\problem{
或問乙從東門直行一十六步，甲從乾隅東行三百二十步望乙。問答同前。
\method{甲東行內減二之乙東行，復以乘甲東行為實，從空，四之甲東行
內減二之乙東行為廉，四益隅，得城徑。}
}

\setcounter{probnum}{16}
\problem{
或問乙從東門直行一十六步，甲從乾隅東行三百二十步望乙，與城參相
直。問答同前。
\method{甲東行內減二之乙東行，復以乘甲東行為實，從空，四之甲東行
內減二之乙東行為廉，四益隅，得城徑。}
}

\setcounter{probnum}{17}
\problem{
或問乙從東門直行一十六步，甲從乾隅東行三百二十步望乙。問答同前。
\method{甲東行內減二之乙東行，復以乘甲東行為實，從空，四之甲東行
內減二之乙東行為廉，四益隅，得城徑。}
}

\setcounter{probnum}{18}
\problem{
或問乙從東門直行一十六步，甲從乾隅東行三百二十步望乙，與城參相
直。問答同前。
\method{甲東行內減二之乙東行，復以乘甲東行為實，從空，四之甲東行
內減二之乙東行為廉，四益隅，得城徑。}
}

\vspace{0.5cm}
\longline
\begin{center}
\textit{測圓海鏡卷六終}
\end{center}

\newpage

%% ===================================================================
%% APPENDIX: STRUCTURE OF THE WORK
%% ===================================================================
\section*{附錄：全書結構}
\addcontentsline{toc}{section}{附錄：全書結構}

《測圓海鏡》全書十二卷，共一百七十問：

\begin{center}
\begin{tabular}{clc}
\toprule
\textbf{卷} & \textbf{篇名} & \textbf{問數} \\
\midrule
卷一 & 圓城圖式 & --- \\
卷二 & 正率一十四問 & 14 \\
卷三 & 邊股一十七問 & 17 \\
\rowcolor{gray!15}
卷四 & 底勾一十七問 & 17 \\
\rowcolor{gray!15}
卷五 & 大股一十八問 & 18 \\
\rowcolor{gray!15}
卷六 & 大勾一十八問 & 18 \\
卷七 & 明勾一十六問 & 16 \\
卷八 & 明股一十六問 & 16 \\
卷九 & 雜題 & 18 \\
卷十 & 雜題 & 18 \\
卷十一 & 雜題 & 18 \\
卷十二 & 雜題 & 10 \\
\midrule
 & \textbf{合計} & \textbf{170} \\
\bottomrule
\end{tabular}
\end{center}

\vspace{0.5cm}
\noindent\textbf{術語解釋：}
\begin{itemize}
  \item \textbf{天元一}：設未知數為一，即現代代數之$x$
  \item \textbf{寄左}：將某式暫存一邊，相當於方程之一邊
  \item \textbf{相消}：兩式相減，相當於移項合併
  \item \textbf{帶分母}：某式含有分母未除
  \item \textbf{冪}：平方，即現代記號$x^2$
  \item \textbf{開方}：開平方，即求平方根
  \item \textbf{開立方}：開立方，即求立方根
  \item \textbf{實}：常數項
  \item \textbf{從}：一次項係數
  \item \textbf{廉}：二次項係數
  \item \textbf{隅}：三次項係數（立方方程中）
\end{itemize}

\vspace{1cm}
\begin{center}
\rule{0.3\textwidth}{0.4pt}

\textit{--- 全書完 ---}
\end{center}

\end{document}
